\begin{table}[!htbp]
\centering
\caption{Expected value added by applicant continuation group}
\label{tab:expected-va-by-continuation}
\resizebox{\textwidth}{!}{%
\begin{tabular}{lccccccc}
\toprule
 & \multicolumn{2}{c}{Achievement} & \multicolumn{3}{c}{Postsecondary access} & \multicolumn{2}{c}{Within higher-education choice} \\
\cmidrule(lr){2-3} \cmidrule(lr){4-6} \cmidrule(lr){7-8}
 & Math & Verbal & Exam taking & Fin. aid app. & Higher-ed enrollment & High-premium inst. & High-premium field \\
\midrule
School ending & -0.108 & -0.087 & -4.652 & 0.522 & -1.126 & -3.246 & 0.433 \\
Voluntary & -0.089 & -0.065 & -1.233 & 1.518 & -0.173 & -2.460 & 0.413 \\
\bottomrule
\end{tabular}
}
\par\medskip\footnotesize
\begin{minipage}{\textwidth}
\textit{Notes:} Assignment-risk fractions use all applicants in each Table 2 subgroup before complete-control and outcome restrictions. School ending means the grade-8 origin has no SAE offering above grade 8; voluntary means it offers grade 9. These operational labels describe continuation availability, not observed motives or actual moves. Expected VA is the mean of sum_j p_ij times the original EB VA_j, without recentering within either group. The main expected-VA table uses all subgroup applicants; the sensitivity uses the exact outcome-specific IV estimation samples. Math and verbal expected VA use outcome-SD units. Binary-outcome expected VA is multiplied by 100 and expressed in percentage points. Pass-through and first-stage coefficients are not rescaled. Expected VA reflects portfolios, priorities and assignment probabilities together; it is not a pure measure of ranked preferences. As in the main paper, unmatched outcomes and school-options without measured VA contribute zero, and probability mass is not renormalized. The coverage table shows the mean assignment probability on schools with measured VA, so differences in coverage can be assessed.
\end{minipage}
\end{table}
